This paper presents the design, implementation, and simulation-based validation of a thrust vector
controlled model rocket. The system uses a custom two-axis gimbal to deflect motor thrust in real time,
replacing passive aerodynamic fins as the sole means of attitude stabilization. A closed-loop PID controller continuously corrects pitch and yaw deviatiohns, while a Kalman filter fuses gyroscope and accelerometer data to produce reliable orientation estimates under vibration and high-acceleration conditions. 
The flight software is structured as an autonomous state machine, managing all phases of
flight from launch detection through apogee identification and recovery. Because powered flight testing
was not possible due to airspace restrictions, system performance was evaluated entirely through
physics-based simulation, including disturbance rejection analysis and a two-axis Monte Carlo study.
These simulations validate the control architecture, achieving a 100\% recovery rate across a two-axis Monte Carlo study and defining a maximum recoverable tilt of approximately 12\degree{}, the physical boundary of stable flight.
The broader aim of this work is to demonstrate that advanced control techniques like thrust vectoring
are achievable at the hobby and educational scale, and to provide a foundation that other students can
build on.
